% Section V template.  tools/build_section.py fills every double-brace placeholder from the logs, CSV tables and configuration of the
% final run (results/section_numbers.json) and writes results/experimental_section.tex; no number in the prose is typed
% by hand.  To be \input from the main paper (macros and labels of Sections II-IV and Appendices A-D); paths are
% relative to results/.  Packages: graphicx, booktabs, amssymb.
% Review round: every figure goes through \pgcfig, which keeps the natural size of the PDF (3.5in single panel, 4.4in
% multi-panel) and caps it at width <= min(\linewidth, 0.62\textwidth) and height <= 0.72\textheight (keepaspectratio),
% so the same source fits the one-column and the two-column IEEEtran layout; the final size is written to the log.
\makeatletter
\@ifundefined{pgcfigw}{\newlength{\pgcfigw}\newsavebox{\pgcfigbox}}{}
\makeatother
\providecommand{\pgcfig}[1]{\sbox{\pgcfigbox}{\includegraphics{#1}}\setlength{\pgcfigw}{\wd\pgcfigbox}\ifdim\pgcfigw>0.62\textwidth\setlength{\pgcfigw}{0.62\textwidth}\fi\ifdim\pgcfigw>\linewidth\setlength{\pgcfigw}{\linewidth}\fi\sbox{\pgcfigbox}{\includegraphics[width=\pgcfigw,height=0.72\textheight,keepaspectratio]{#1}}\typeout{PGCFIG #1 wd=\the\wd\pgcfigbox\space ht=\the\ht\pgcfigbox\space tw=\the\textwidth\space th=\the\textheight\space lw=\the\linewidth}\usebox{\pgcfigbox}}
\section{Numerical Validation}
\label{sec:numerics}

\subsection{Benchmark, implementation and protocol}
The benchmark is the class-(H) telegraph reaction--diffusion--advection lattice
\begin{align*}
\dot q &= p,\\
\dot p &= -(\nu L_h + \alpha I)q - \gamma p - c\,W_f D_h W_f\, q\\
&\quad - \kappa(q - \sin q) + Bu + B_w w,
\end{align*}
on \(N_m = {{N_m}}\) nodes (\(N = {{N}}\) states) with Dirichlet data, \(\nu = {{nu}}\), \(\alpha = {{alpha}}\),
\(\gamma = {{gamma}}\), \(\kappa = {{kappa}}\), two actuators (\(m = {{m}}\)) and one disturbance channel. Here \(D_h\) is the
skew central difference and the advection filter is
\[
W_f = \Phi_f \operatorname{diag}(\underbrace{1,\dots,1}_{{{n_adv_full_weight}}},\,w,\dots,w)\,\Phi_f^\top ,
\]
where \(\Phi_f\) holds the first \(K_f = {{n_adv_modes}}\) mass-normalised sine modes. The plain projector
\(\Pi_f = \Phi_f\Phi_f^\top\) is the special case \(w = {{e9_wmax}}\). The cutoff-mode weight \(w\) is a disclosed benchmark
parameter: it is a configuration literal, it is swept in Experiment 9, and the deployed value
\(w^\star = {{w_star_exp}} \approx {{w_star}}\) is the one selected by that sweep. The ROM keeps \(K_m = {{K_m}}\)
eigenmodes (\(K = {{K}}\)), so \(K_f > K_m\) and the advection couples the ROM to the tail.

The advection is not folded into the damped-wave operator. \(A_{\mathrm{full}} = (J_{\mathrm{full}} - R_{\mathrm{full}})P_{\mathrm{full}}\)
stays the pure damped wave, so \(Q = \gamma P\) holds (relative error \({{Q_identity_relerr_R2}}\)), \(\lambda_{K+1} = {{lambda_K1_R2}}\)
is unchanged, and (H2a) holds by eigenmode invariance (Lemma~\ref{lem:h2a_char}, residual \({{H2a_residual_R2}}\)). For
the same reason the block \(P_0^{-1}\Phi^\top P_{\mathrm{full}}A_{\mathrm{full}}P_\perp\) vanishes to round-off
(\({{A_rperp_full_norm_R2}}\)). The linear cross-coupling sits in the nonlinearity,
\(F_{\mathrm{nl}}(x) = A_{\mathrm{adv}}x + F_{\sin}(x)\) with
\[
A_{\mathrm{adv}} = \begin{bmatrix}0&0\\-c\,W_fD_hW_f&0\end{bmatrix},
\]
and its ROM--tail block \(A_{\mathrm{adv},r,\perp} := P_0^{-1}\Phi^\top P_{\mathrm{full}}A_{\mathrm{adv}}P_\perp\) of
Lemma~\ref{lem:delta_cert_classH} has norm \(\|A_{\mathrm{adv},r,\perp}\|_2 = {{A_adv_rperp_norm_R2}}\) in R2 and
\({{A_adv_rperp_norm_R1}}\) in R1.

Two regimes are used. R2 has \(c = {{c_R2}}\) (\(\mathrm{Pe} = {{Pe_R2}}\)) and is open-loop unstable. R1 has
\(c = {{c_R1}}\) (\(\mathrm{Pe} = {{Pe_R1}}\)). The sampling period of R2 is \(T_s = {{Ts_R2}}\) with
\(\tau_A = {{tauA_R2}}\), \(\tau_B = {{tauB_R2}}\). This is a trade-off forced by the active tail coupling: at the
previous \(T_s = {{Ts_wave2}}\) the certificate refuses at \(w^\star\), because
\(\delta_{\mathrm{total}} = {{wave2_delta_total}} > \delta_{\mathrm{cert}} = {{wave2_delta_cert}}\) and
the authority condition (H6), \eqref{eq:control_authority}, fails (Table~\ref{tab:exp1_refusals}). R1 keeps \(T_s = {{Ts_R1}}\).
The short period is set by the hold term, not by the plant time scale: the slowest mode (\(\omega_{\min} = {{omega_min_R2}}\))
is sampled \({{samples_per_period_R2}}\) times per period \(2\pi/\omega_{\min}\) in R2 (\({{samples_per_period_R1}}\) in R1),
and the hold term \(\phi L_{\mathrm{hold}}T_s\) of \eqref{eq:delta_total_exact}, with \(L_{\mathrm{hold}} = {{L_hold_R2}}\)
and \(\mathrm{Lip}(L_fh) = {{L_Lfh_R2}}\), is a share \({{hold_share_R2}}\) of \(\delta_{\mathrm{total}}\) in R2
(\({{hold_share_R1}}\) in R1, where \(L_{\mathrm{hold}} = {{L_hold_R1}}\) and \(\mathrm{Lip}(L_fh) = {{L_Lfh_R1}}\)).

Every constant of Sections~\ref{sec:problem}--\ref{sec:theory} is computed at run time in float64 by a function that
names the equation it implements. A pytest AST scan fails if a float literal appears in the numerical modules, and a
perturbation test fails if any certificate quantity does not depend on the configuration inputs. The certificate raises
(refuses) unless (H0)--(H7), the timing order and the codesign window all hold. This single flag is the
Theorem~\ref{thm:main_cbf} hypothesis row of Table~\ref{tab:exp1_hypotheses}. (H-dz) is checked separately and only
for Proposition~\ref{prop:deadzone_stability}. The Path-A LUT is indexed in port space \(v = g_r^\top P z\) by default.
When \(\mathcal C_{r,0} = \{\|z\|_P^2 \le E_{\mathrm{safe}}\}\) as in \eqref{eq:Cr0}, its range is
\(D_v = 2\sqrt{E_{\mathrm{safe}}}\max_j \sqrt{(g_r^\top P g_r)_{jj}} = {{D_v_R2}}\) by \eqref{eq:Dv_closed}; the largest
sampled value is \({{D_v_sampled_R2}}\) (both to four significant digits).
The modulus \(\mathrm{Lip}(L_fh)\) that enters \(L_{\mathrm{hold}}\) carries the factor of two of the gradient
\(\nabla_z(z^\top P z) = 2Pz\); it is cross-checked against finite-difference quotients of \(L_fh\) at random pairs in
\(\mathcal C_{r,0}\), which must not exceed the certified modulus.

The full-order plant is integrated with an exponential integrator (ETD2RK in the lattice eigenbasis). The latch
instants \(\tau_A\), \(\tau_B\) are integration nodes, so the switching times are exact. The disturbance is the worst
case of Remark~\ref{rem:hji_game} with \(\epsilon_{\mathrm{adv}} = {{eps_adv}}\). All seeds come from the configuration
file. The suite runs {{n_assert}} experiment assertions and {{n_tests}} pytest tests on two CPU cores in
\({{rt_total}}\)~s, with a peak resident memory of \({{rt_peak}}\)~MB (Table~\ref{tab:runtimes}).

\begin{table*}[t]
\centering
\caption{Certificate constants (part 1 of 2), computed at run time in float64 from the lattice operators. $K_f$ is the advection band, $w$ the disclosed cutoff-mode weight (Table~\ref{tab:exp9_wsweep}); $L_{\perp,P}(\mathrm{R2})/L_{\perp,P}(\mathrm{R1})=14.82$.}
\label{tab:exp1_constants}
\footnotesize
\setlength{\tabcolsep}{4pt}
\begin{tabular}{lrrr}
\toprule
quantity & equation & R2 (Pe $= 20$) & R1 (Pe $= 1$) \\
\midrule
$N_m$ & -- & $160$ & $160$ \\
$N$ & -- & $320$ & $320$ \\
$m$ & -- & $2$ & $2$ \\
$\nu$ & -- & $0.05$ & $0.05$ \\
$\alpha$ (plant) & -- & $1$ & $1$ \\
$\gamma$ & -- & $1$ & $1$ \\
$\kappa$ & -- & $0.05$ & $0.05$ \\
$c$ & -- & $1$ & $0.05$ \\
$K_m$ & -- & $2$ & $2$ \\
$K$ & -- & $4$ & $4$ \\
full-weight advection modes & -- & $2$ & $2$ \\
$E_{\mathrm{safe}}$ & \eqref{eq:H0} & $0.05$ & $0.05$ \\
$\varepsilon_{\mathrm{adv}}$ & \eqref{eq:M_perp} & $0.002$ & $0.002$ \\
open-loop growth $\max\operatorname{Re}\lambda(A_{\mathrm{lin}})$ & -- & $0.2925$ & $-0.5$ \\
$K_f$ & -- & $4$ & $4$ \\
$w$ (cutoff weight) & -- & $0.0221$ & $0.0221$ \\
$\mathrm{Pe}$ & Remark~\ref{rem:rda_cutoff} & $20$ & $1$ \\
$\lambda_{\min}(P_{\mathrm{full},\epsilon})$ & \eqref{eq:eps_choice} & $0.6892$ & $0.6892$ \\
$\omega_{\min}$ & \eqref{eq:eps_choice} & $1.222$ & $1.222$ \\
$\mu_{2,P}(A_{r,0})$ & \eqref{eq:mu_pe} & $-0.5$ & $-0.5$ \\
$\|Q-\gamma P\|/\|\gamma P\|$ & \eqref{eq:mu_pe} & $7.039\times10^{-15}$ & $7.039\times10^{-15}$ \\
$\lambda_{K+1}$ & \eqref{eq:tail_dissipation} & $0.5$ & $0.5$ \\
(H2a) residual $\|\Phi^\top P_{\mathrm{full},\epsilon}P_\perp\|$ & Lemma~\ref{lem:h2a_char} & $6.348\times10^{-13}$ & $6.348\times10^{-13}$ \\
$\|A_{r,\perp}\|_2$ (eigenmode invariance) & Lemma~\ref{lem:h2a_char} & $1.27\times10^{-12}$ & $1.27\times10^{-12}$ \\
$\rho_{\mathrm{ODD}}$ & \eqref{eq:H0} & $0.2828$ & $0.2828$ \\
$D_{\mathrm{ODD}}$ & \eqref{eq:Cr0} & $0.5656$ & $0.5656$ \\
$\mathrm{Lip}(F_{\mathrm{nl}})$ & Assumption~\ref{assump:generic} & $2.67$ & $0.1354$ \\
$M_\perp$ & \eqref{eq:M_perp} & $0.04916$ & $0.003647$ \\
$\varepsilon_{\mathrm{basis},P}$ & \eqref{eq:eps_basis_P} & $0.09831$ & $0.007293$ \\
$\|A_{\mathrm{adv},r,\perp}\|_2$ & Lemma~\ref{lem:delta_cert_classH} & $0.106$ & $0.005301$ \\
$L_{\perp,P}$ (advective, linear) & \eqref{eq:Lperp_def} & $0.04654$ & $0.002327$ \\
$L_{\perp,P}$ ($\sin$ nonlinearity) & \eqref{eq:Lperp_def} & $8.719\times10^{-4}$ & $8.719\times10^{-4}$ \\
$L_{\perp,P}$ & \eqref{eq:Lperp_def} & $0.04741$ & $0.003199$ \\
linear share of $L_{\perp,P}$ & \eqref{eq:Lperp_def} & $0.9816$ & $0.7274$ \\
$L_h$ & \eqref{eq:reduced_constants} & $0.7865$ & $0.7865$ \\
$g_h$ & \eqref{eq:reduced_constants} & $0.4788$ & $0.4788$ \\
$d_h$ & \eqref{eq:reduced_constants} & $0.4472$ & $0.4472$ \\
$M_z$ & \eqref{eq:reduced_constants} & $4.324$ & $0.7887$ \\
$\mathrm{Lip}(L_fh)$ & \eqref{eq:reduced_constants} & $1.77$ & $0.7948$ \\
$\mathrm{Lip}(L_gh)$ & \eqref{eq:reduced_constants} & $2.394$ & $2.394$ \\
$L_{\mathcal{T},z}$ & \eqref{eq:LTz} & $5.985$ & $0.5985$ \\
$\eta_{\mathcal{T}}$ & Lemma~\ref{lem:sector} & $4.975$ & $0.4975$ \\
(H-sat) left side $L_{\mathcal{T},z}\sqrt{K}D_{\mathrm{ODD}}/2$ & \eqref{eq:sat_cond} & $3.386$ & $0.3386$ \\
$u_{\max}$ & \eqref{eq:sat_cond} & $3.4$ & $0.35$ \\
$\varepsilon_{\mathrm{CROWN}}$ (backward) & \eqref{eq:twin_crown_bound} & $0.004455$ & $4.536\times10^{-4}$ \\
\bottomrule
\end{tabular}
\end{table*}

\begin{table*}[t]
\centering
\caption{Certificate constants (part 2 of 2).}
\label{tab:exp1_constants_b}
\footnotesize
\setlength{\tabcolsep}{4pt}
\begin{tabular}{lrrr}
\toprule
quantity & equation & R2 (Pe $= 20$) & R1 (Pe $= 1$) \\
\midrule
$E_q$ (explicit) & \eqref{eq:crown_explicit} & $0.005674$ & $5.58\times10^{-4}$ \\
$\varepsilon_{\mathrm{CROWN}}$ & Prop.~\ref{prop:twin_crown} & $0.004455$ & $4.536\times10^{-4}$ \\
$D_v$ & \eqref{eq:Dv_closed} & $0.4472$ & $0.4472$ \\
$D_v$ (sampled sup) & \eqref{eq:Dv_closed} & $0.4472$ & $0.4472$ \\
$M_v$ & \eqref{eq:eps_A_triangle} & $5.176$ & $0.9441$ \\
$\Delta_{q,A}$ (LUT step) & \eqref{eq:Dv_closed} & $8.735\times10^{-4}$ & $8.735\times10^{-4}$ \\
$\varepsilon_A$ & \eqref{eq:eps_A_triangle} & $0.007543$ & $7.624\times10^{-4}$ \\
$\Delta u_{\mathrm{lag}}$ & \eqref{eq:eps_A_triangle} & $0.007765$ & $9.441\times10^{-4}$ \\
window lower bound & \eqref{eq:codesign_window} & $8$ & $8$ \\
$b_A$ & \eqref{eq:codesign_window} & $9$ & $9$ \\
window upper bound & \eqref{eq:codesign_window} & $11$ & $11$ \\
$T_s$ & \eqref{eq:multirate_hold} & $3\times10^{-4}$ & $0.002$ \\
$\tau_A$ & \eqref{eq:multirate_hold} & $1.5\times10^{-5}$ & $1\times10^{-4}$ \\
$\tau_B$ & \eqref{eq:multirate_hold} & $3\times10^{-5}$ & $2\times10^{-4}$ \\
samples per period $2\pi/(\omega_{\min}T_s)$ & \eqref{eq:multirate_hold} & $1.714\times10^{4}$ & $2571$ \\
$b_B$ & Prop.~\ref{prop:twin_crown} & $12$ & $12$ \\
$\gamma_0$ & \eqref{eq:teacher_cbf} & $0.35$ & $0.3$ \\
$w_1$ & \eqref{eq:exp_weights} & $0.05$ & $0.04999$ \\
$w_2$ & \eqref{eq:exp_weights} & $0.05$ & $0.04999$ \\
$w_3$ & \eqref{eq:exp_weights} & $0.9$ & $0.9$ \\
$\phi(\gamma_0T_s)$ & \eqref{eq:delta_total_exact} & $0.5$ & $0.5$ \\
$L_{\mathrm{hold}}$ & \eqref{eq:delta_dip} & $44.04$ & $1.474$ \\
$\Delta h_{\mathrm{dip}}$ & \eqref{eq:delta_dip} & $2.06\times10^{-6}$ & $3.008\times10^{-6}$ \\
$\delta_{\mathrm{total}}$: tail term & \eqref{eq:delta_total_exact} & $0.003666$ & $1.835\times10^{-5}$ \\
$\delta_{\mathrm{total}}$: quantisation term & \eqref{eq:delta_total_exact} & $0.002207$ & $2.246\times10^{-4}$ \\
$\delta_{\mathrm{total}}$: lag term & \eqref{eq:delta_total_exact} & $1.859\times10^{-4}$ & $2.26\times10^{-5}$ \\
$\delta_{\mathrm{total}}$: adversary term & \eqref{eq:delta_total_exact} & $8.944\times10^{-4}$ & $8.944\times10^{-4}$ \\
$\delta_{\mathrm{total}}$: hold term & \eqref{eq:delta_total_exact} & $0.006607$ & $0.001474$ \\
$\delta_{\mathrm{total}}$ & \eqref{eq:delta_total_exact} & $0.01356$ & $0.002634$ \\
tail share $L_hL_{\perp,P}\varepsilon_{\mathrm{basis},P}/\delta_{\mathrm{total}}$ & \eqref{eq:delta_total_exact} & $0.2704$ & $0.006967$ \\
hold share of $\delta_{\mathrm{total}}$ & \eqref{eq:delta_total_exact} & $0.4872$ & $0.5596$ \\
$h_{\mathrm{act}}$ & \eqref{eq:shell} & $0.02068$ & $0.02559$ \\
$r_{\min}$ & \eqref{eq:shell} & $0.1402$ & $0.156$ \\
$\lambda_{\min}(M,P)$ & Lemma~\ref{lem:delta_cert_classH} & $0.3731$ & $0.9835$ \\
$s_{\mathrm{nl}}$ & Lemma~\ref{lem:delta_cert_classH} & $3.005\times10^{-4}$ & $3.005\times10^{-4}$ \\
$c_M$ & Lemma~\ref{lem:delta_cert_classH} & $0.3725$ & $0.9829$ \\
$\delta_{\mathrm{cert}}$ & Lemma~\ref{lem:delta_cert_classH} & $0.01456$ & $0.03161$ \\
$\delta_{\mathrm{cert}}$ branch & \eqref{eq:delta_cert_classH} & $1$ & $1$ \\
margin $\delta_{\mathrm{cert}}-\delta_{\mathrm{total}}$ & \eqref{eq:control_authority} & $9.994\times10^{-4}$ & $0.02898$ \\
$\gamma_0h_{\mathrm{act}}$ (Cor.~\ref{cor:stochastic_w2}) & Cor.~\ref{cor:stochastic_w2} & $0.007237$ & $0.007678$ \\
null-space rate & \eqref{eq:null_space_cond} & $1$ & $1$ \\
$\mu_R$ & Prop.~\ref{prop:deadzone_stability} & $0.5$ & $0.5$ \\
$L_{\mathrm{sec}}$ & Prop.~\ref{prop:deadzone_stability} & $2.392$ & $0.1204$ \\
$\bar\varepsilon_q$ & Prop.~\ref{prop:deadzone_stability} & $0.01998$ & $0.002342$ \\
$r_{\mathrm{ss}}$ & \eqref{eq:rss_ball} & $\infty$ & $0.01162$ \\
\bottomrule
\end{tabular}
\end{table*}

\begin{table*}[t]
\centering
\caption{Hypothesis ledger of Theorem~\ref{thm:main_cbf}, hypotheses (H0)--(H7), and of Proposition~\ref{prop:deadzone_stability}, conditions (i) $\mu_R>L_{\mathrm{sec}}$ and (H-dz); \checkmark{} = satisfied.}
\label{tab:exp1_hypotheses}
\footnotesize
\setlength{\tabcolsep}{4pt}
\begin{tabular}{llcc}
\toprule
hypothesis & content & R2 & R1 \\
\midrule
(H0) & ODD containment and $E_{\mathrm{safe}}$ & \checkmark & \checkmark \\
(H1) & Assumption~\ref{assump:generic} constants on $\mathcal{X}_{\mathrm{ODD}}$ & \checkmark & \checkmark \\
(H2) & (H2a)--(H2c), $\gamma<2\omega_{\min}$, $\epsilon=\gamma/2$ & \checkmark & \checkmark \\
(H3) & Teacher: $H/2\ge m$, $G_{\mathrm{port}}\succ0$, (H-sat) & \checkmark & \checkmark \\
(H4) & tail tolerance, $h_{\mathrm{act}}$ interval & \checkmark & \checkmark \\
(H5) & $b_A\ge1$ satisfies \eqref{eq:lut_capacity} & \checkmark & \checkmark \\
(H6) & authority \eqref{eq:control_authority} & \checkmark & \checkmark \\
(H7) & initial conditions & \checkmark & \checkmark \\
(H-sat) & $L_{\mathcal{T},z}\sqrt{K}D_{\mathrm{ODD}}/2\le u_{\max}$ \eqref{eq:sat_cond} & \checkmark & \checkmark \\
timing $0<\tau_A<\tau_B<T_s$ & latch order of \eqref{eq:multirate_hold} & \checkmark & \checkmark \\
codesign window \eqref{eq:codesign_window} & $b_A$ inside the window & \checkmark & \checkmark \\
Thm.~\ref{thm:main_cbf} hypotheses & all of the above & \checkmark & \checkmark \\
$\mu_R>L_{\mathrm{sec}}$ & Prop.~\ref{prop:deadzone_stability}(i) contraction rate & $\times$ & \checkmark \\
(H-dz) & $r_{\mathrm{ss}}<r_{\min}$ (Prop.~\ref{prop:deadzone_stability} only) & $\times$ & \checkmark \\
Prop.~\ref{prop:deadzone_stability} hypotheses & (i) $\mu_R>L_{\mathrm{sec}}$, (H-dz) and (H3) & $\times$ & \checkmark \\
\bottomrule
\end{tabular}
\end{table*}


\subsection{Experiment 0: open-loop exit}
With \(\mathrm{Pe} = {{Pe_R2}}\) and \(w = w^\star\) the linearisation has growth rate \({{open_loop_growth_R2}}\)
(Table~\ref{tab:exp1_constants}). All {{e0_exits}} of
{{n_seeds}} initial conditions satisfying (H7) leave \(\mathcal C\) without control: the worst \(\min_t h\) is
\({{e0_minh}}\) and the latest exit is at \(t = {{e0_lastexit}}\) (Fig.~\ref{fig:exp0}, Table~\ref{tab:exp0_openloop}).
The safety problem is therefore not trivial: invariance in the later experiments comes from the controller, not from
passivity.

\begin{table}[t]
\centering
\caption{Experiment 0: open-loop ($u\equiv0$) trajectories of the R2 benchmark (Pe $= 20$) from certified initial states.}
\label{tab:exp0_openloop}
\footnotesize
\setlength{\tabcolsep}{4pt}
\begin{tabular}{lrrr}
\toprule
seed & $h(x_0)$ & $\min_t h(x(t))$ & first exit time \\
\midrule
$0$ & $2.06\times10^{-6}$ & $-0.7165$ & $3\times10^{-4}$ \\
$1$ & $2.06\times10^{-6}$ & $-0.04703$ & $3.856$ \\
$2$ & $2.06\times10^{-6}$ & $-0.1102$ & $2.983$ \\
$3$ & $2.06\times10^{-6}$ & $-0.6174$ & $0.7137$ \\
$4$ & $2.06\times10^{-6}$ & $-0.5155$ & $0.4689$ \\
$5$ & $2.06\times10^{-6}$ & $-0.6655$ & $0.0657$ \\
$6$ & $2.06\times10^{-6}$ & $-0.6377$ & $9\times10^{-4}$ \\
$7$ & $2.06\times10^{-6}$ & $-0.4675$ & $0.8838$ \\
$8$ & $2.06\times10^{-6}$ & $-0.606$ & $0.4128$ \\
$9$ & $2.06\times10^{-6}$ & $-0.8979$ & $3\times10^{-4}$ \\
$10$ & $2.06\times10^{-6}$ & $-0.2972$ & $1.566$ \\
$11$ & $2.06\times10^{-6}$ & $-0.5958$ & $0.2889$ \\
$12$ & $2.06\times10^{-6}$ & $-0.2189$ & $2.061$ \\
$13$ & $2.06\times10^{-6}$ & $-0.351$ & $1.379$ \\
$14$ & $2.06\times10^{-6}$ & $-0.06081$ & $3.636$ \\
$15$ & $2.06\times10^{-6}$ & $-0.6837$ & $3\times10^{-4}$ \\
\bottomrule
\end{tabular}
\end{table}

\begin{figure}[t]\centering\pgcfig{figures/exp0_openloop_h.pdf}
\caption{Experiment 0: raw open-loop \(h(x(t))\) in R2 (\(\mathrm{Pe}={{Pe_R2}}\)) for all seeds.}\label{fig:exp0}\end{figure}

\subsection{Experiment 1: constants, hypotheses and refusals}
Tables~\ref{tab:exp1_constants} and~\ref{tab:exp1_constants_b} list every certificate constant for R2 and R1, and
Table~\ref{tab:exp1_hypotheses} lists the hypotheses. Proposition~\ref{prop:dahlquist} holds to round-off:
\(\mu_{2,P}(A_{r,0}) = {{mu2P}}\).

The tail coupling is now active. In R2, \(\|A_{\mathrm{adv},r,\perp}\|_2 = {{A_adv_rperp_norm_R2}}\) and
\(L_{\perp,P} = {{L_perp_R2}}\). Of this, \({{L_perp_lin_R2}}\) is the linear advective part and \({{L_perp_nl_R2}}\)
the \(\sin\) nonlinearity, so the advective part is a fraction \({{L_perp_lin_share_R2}}\) of the total. In R1,
\(L_{\perp,P} = {{L_perp_R1}}\), and the ratio between the regimes is \({{Lperp_ratio}}\). The tail term
\(L_h L_{\perp,P}\varepsilon_{\mathrm{basis},P}\) is a share \({{tail_share_R2}}\) of
\(\delta_{\mathrm{total}} = {{delta_total_R2}}\) in R2 (and a share \({{tail_share_R1}}\) of
\(\delta_{\mathrm{total}} = {{delta_total_R1}}\) in R1). With
\(\varepsilon_{\mathrm{basis},P} = {{eps_basis_P_R2}}\), the R2 margin is
\(\delta_{\mathrm{cert}} - \delta_{\mathrm{total}} = {{delta_cert_R2}} - {{delta_total_R2}} = {{H6_margin_R2}}\).

Lemma~\ref{lem:delta_cert_classH} gives \(c_M = {{c_M_R2}}\) from \(\lambda_{\min}(M,P) = {{lambda_min_MP_R2}}\) and
\(s_{\mathrm{nl}} = {{s_nl_R2}}\) in R2, and \(c_M = {{c_M_R1}}\) from \(\lambda_{\min}(M,P) = {{lambda_min_MP_R1}}\) and
\(s_{\mathrm{nl}} = {{s_nl_R1}}\) in R1. The codesign window is
\([{{codesign_lo}}, {{codesign_hi}}]\) and contains \(b_A = {{b_A}}\).

Table~\ref{tab:exp1_refusals} shows the certificate refusing {{n_refusals}} inadmissible variants. The state-space LUT (the
\(m = K\) variant of Remark~\ref{rem:lut_variants}) has an empty window \([{{ss_window}}]\). The other refusals are
\(u_{\max}\) below the (H-sat) left side, unfiltered lattice advection, the plain projector \(\Pi_f\), the previous
sampling period and \(T_s\), \(\tau_A\), \(\tau_B\) scaled by {{Ts_refusal_factor}}. Proposition~\ref{prop:deadzone_stability} needs (i) \(\mu_R > L_{\mathrm{sec}}\), (H-dz) and (H3)
(Table~\ref{tab:exp1_hypotheses}). In R1 all three hold: \(\mu_R = {{mu_R_R1}} > L_{\mathrm{sec}} = {{L_sec_R1}}\) and
\(r_{\mathrm{ss}} = {{r_ss_R1}} < r_{\min} = {{r_min_R1}}\). In R2, \(L_{\mathrm{sec}} = {{L_sec_R2}}\) exceeds
\(\mu_R = {{mu_R_R2}}\), so condition (i) fails and Proposition~\ref{prop:deadzone_stability} does not apply; this does
not affect Theorem~\ref{thm:main_cbf}.

\begin{table*}[t]
\centering
\caption{Certificate refusals on non-admissible variants of R2. The LUT window is \eqref{eq:codesign_window} in the variant's own form (port space by default; the state-space row uses $K$, $D_{\mathrm{ODD}}$ and $M_z$).}
\label{tab:exp1_refusals}
\footnotesize
\setlength{\tabcolsep}{4pt}
\begin{tabular}{lrrrrr}
\toprule
configuration & refused & failing & $\delta_{\mathrm{total}}$ & $\delta_{\mathrm{cert}}$ & window $[\mathrm{lo},\mathrm{hi}]$ \\
\midrule
state-space LUT ($m = K$ variant) & \checkmark & (H5), codesign window \eqref{eq:codesign_window} & $0.01377$ & $0.01456$ & [9, 5] \\
$u_{\max}$ below the (H-sat) left side & \checkmark & (H3) & $0.01245$ & $0.01483$ & [8, 11] \\
pure lattice advection (no spectral cutoff) & \checkmark & (H4), (H6), (H7) & $157.6$ & $-3091$ & [3, 11] \\
plain projector $\Pi_f$ ($w = 1$) & \checkmark & (H4), (H6), (H7) & $7.93$ & $-8.456$ & [6, 11] \\
previous sampling period $T_s=5\times10^{-4}$ at $w^\star$ & \checkmark & (H6) & $0.01821$ & $0.01455$ & [7, 11] \\
$T_s$, $\tau_A$, $\tau_B$ $\times 10$ (hold term) & \checkmark & (H6) & $0.07483$ & $0.01438$ & [5, 11] \\
\bottomrule
\end{tabular}
\end{table*}


\subsection{Experiment 9: the cutoff-mode weight and the certifiable frontier}
The weight \(w\) is swept over {{e9_n}} geometric values from \({{e9_wmax}}\) down to \({{e9_wmin}}\) at the R2 P\'eclet number
(Table~\ref{tab:exp9_wsweep}, Fig.~\ref{fig:exp9}). At every \(w\) the open loop leaves \(\mathcal C\) on all seeds, so
the plant is unsafe throughout the sweep. For \(w \ge {{e9_wH4}}\), the tail tolerance (H4) fails together with (H6) and
(H7). For \({{e9_H6_lo}} \le w \le {{e9_H6_hi}}\), only the authority condition (H6) fails. The certificate accepts
for \(w \le w^\star = {{e9_wstar}}\). Below \(w^\star\) the single-rate ablation (c) is no longer refused (first at
\(w = {{e9_c_acc}}\)). The deployed \(w^\star\) is therefore the largest grid value that is accepted with a positive
authority margin \(\delta_{\mathrm{cert}}-\delta_{\mathrm{total}}\) in \eqref{eq:control_authority} while every ablation still breaks. \(\|A_{\mathrm{adv},r,\perp}\|_2\) is proportional to
\(w\), and the tail share of \(\delta_{\mathrm{total}}\) falls from \({{e9_share_w1}}\) at \(w = {{e9_wmax}}\) to
\({{e9_share_star}}\) at \(w^\star\).

For the plain projector (\(w = {{e9_wmax}}\)), the P\'eclet scan of Table~\ref{tab:exp9_plain_peclet} shows that the certifiable
advection speeds (\(c \le {{e9_plain_cert}}\), \(\mathrm{Pe} \le {{e9_plain_cert_Pe}}\)) and the open-loop-unstable speeds
(\(c \ge {{e9_plain_unst}}\), \(\mathrm{Pe} \ge {{e9_plain_unst_Pe}}\)) are disjoint. Under the Euclidean-ODD supremum that
defines \(M_\perp\) in~\eqref{eq:M_perp}, full-strength cutoff-mode advection is not certifiable at any P\'eclet number
where the open loop is unstable: no softening. The plain-projector result is the \(L^2\)-ODD conservatism of
Remark~\ref{rem:rda_cutoff}. Keeping \(A_{\mathrm{adv}}\) inside \(F_{\mathrm{nl}}\) preserves (H2a) exactly (eigenmode
invariance, Lemma~\ref{lem:h2a_char}) at the price of charging \(A_{\mathrm{adv},r,\perp}\) in \(M_\perp\) and
\(L_{\perp,P}\). The weight \(w\) attenuates the cutoff modes, which is what an \(H^1\)-bounded ODD would do for a smooth
field; certifying \(w = {{e9_wmax}}\) at \(\mathrm{Pe} \ge {{e9_plain_unst_Pe}}\) would require that \(H^1\) formulation,
which is outside the present theorem.

\begin{table*}[t]
\centering
\caption{Disclosed sweep of the cutoff-mode advection weight $w$ at the R2 P\'eclet number ($w=1$: plain projector $\Pi_f$). The deployed $w^\star$ is the largest grid value accepted with a positive authority margin in \eqref{eq:control_authority} and all ablations refused. Tail share: $L_hL_{\perp,P}\varepsilon_{\mathrm{basis},P}/\delta_{\mathrm{total}}$; margin: $\delta_{\mathrm{cert}}-\delta_{\mathrm{total}}$; exits: open-loop exits of $16$ seeds at $c=1$; --: ablations not evaluated (the full certificate already refuses).}
\label{tab:exp9_wsweep}
\footnotesize
\setlength{\tabcolsep}{3pt}
\begin{tabular}{lrrrrrrrr}
\toprule
$w$ & Thm.~\ref{thm:main_cbf} hypotheses & failing & $\varepsilon_{\mathrm{basis},P}$ & tail share & $\|A_{\mathrm{adv},r,\perp}\|_2$ & margin & abl.\ refused & exits \\
\midrule
$1$ & $\times$ & (H4), (H6), (H7) & $4.77$ & $0.9968$ & $4.798$ & $-16.39$ & -- & 16/16 \\
$0.7071$ & $\times$ & (H4), (H6), (H7) & $2.766$ & $0.995$ & $3.393$ & $-6.088$ & -- & 16/16 \\
$0.5$ & $\times$ & (H4), (H6), (H7) & $1.706$ & $0.9911$ & $2.399$ & $-2.492$ & -- & 16/16 \\
$0.3536$ & $\times$ & (H4), (H6), (H7) & $1.113$ & $0.9831$ & $1.696$ & $-1.107$ & -- & 16/16 \\
$0.25$ & $\times$ & (H4), (H6), (H7) & $0.7587$ & $0.9677$ & $1.2$ & $-0.5211$ & -- & 16/16 \\
$0.1768$ & $\times$ & (H4), (H6), (H7) & $0.5331$ & $0.9388$ & $0.8482$ & $-0.2539$ & -- & 16/16 \\
$0.125$ & $\times$ & (H4), (H6), (H7) & $0.3829$ & $0.8879$ & $0.5998$ & $-0.1256$ & -- & 16/16 \\
$0.08839$ & $\times$ & (H4), (H6), (H7) & $0.2801$ & $0.8052$ & $0.4241$ & $-0.06177$ & -- & 16/16 \\
$0.0625$ & $\times$ & (H6) & $0.2086$ & $0.6865$ & $0.2999$ & $-0.02934$ & -- & 16/16 \\
$0.04419$ & $\times$ & (H6) & $0.1585$ & $0.5417$ & $0.212$ & $-0.01265$ & -- & 16/16 \\
$0.03125$ & $\times$ & (H6) & $0.1232$ & $0.395$ & $0.1499$ & $-0.003795$ & -- & 16/16 \\
$0.0221$ & \checkmark & none & $0.09831$ & $0.2704$ & $0.106$ & $9.994\times10^{-4}$ & \checkmark & 16/16 \\
$0.01562$ & \checkmark & none & $0.08073$ & $0.1782$ & $0.07497$ & $0.003661$ & $\times$ & 16/16 \\
$0.01105$ & \checkmark & none & $0.0683$ & $0.1159$ & $0.05301$ & $0.00518$ & $\times$ & 16/16 \\
$0.007812$ & \checkmark & none & $0.05952$ & $0.0758$ & $0.03748$ & $0.006072$ & $\times$ & 16/16 \\
$0.005524$ & \checkmark & none & $0.0533$ & $0.05036$ & $0.02651$ & $0.006613$ & $\times$ & 16/16 \\
$0.003906$ & \checkmark & none & $0.04891$ & $0.03419$ & $0.01874$ & $0.006949$ & $\times$ & 16/16 \\
\bottomrule
\end{tabular}
\end{table*}

\begin{table*}[t]
\centering
\caption{Plain projector $\Pi_f$ ($w=1$): certificate acceptance versus open-loop instability over the advection speed $c$. The two sets are disjoint.}
\label{tab:exp9_plain_peclet}
\footnotesize
\setlength{\tabcolsep}{4pt}
\begin{tabular}{lrrrrrr}
\toprule
$c$ & $\mathrm{Pe}$ & Thm.~\ref{thm:main_cbf} hypotheses & failing & $\varepsilon_{\mathrm{basis},P}$ & tail share $L_hL_{\perp,P}\varepsilon_{\mathrm{basis},P}/\delta_{\mathrm{total}}$ & max Re eig$(A_{\mathrm{full}}+A_{\mathrm{adv}})$ \\
\midrule
$0.005$ & $0.1$ & \checkmark & none & $0.06197$ & $0.05895$ & $-0.5$ \\
$0.01$ & $0.2$ & \checkmark & none & $0.08563$ & $0.1427$ & $-0.5$ \\
$0.02$ & $0.4$ & \checkmark & none & $0.1329$ & $0.3361$ & $-0.5$ \\
$0.03$ & $0.6$ & $\times$ & (H6) & $0.1803$ & $0.5053$ & $-0.5$ \\
$0.05$ & $1$ & $\times$ & (H4), (H6), (H7) & $0.2749$ & $0.7202$ & $-0.5$ \\
$0.1$ & $2$ & $\times$ & (H4), (H6), (H7) & $0.5114$ & $0.9039$ & $-0.5$ \\
$0.2$ & $4$ & $\times$ & (H4), (H6), (H7) & $0.9846$ & $0.9716$ & $-0.5$ \\
$0.3$ & $6$ & $\times$ & (H4), (H6), (H7) & $1.458$ & $0.9859$ & $-0.365$ \\
$0.5$ & $10$ & $\times$ & (H4), (H6), (H7) & $2.404$ & $0.9935$ & $0.1025$ \\
$1$ & $20$ & $\times$ & (H4), (H6), (H7) & $4.77$ & $0.9968$ & $0.8887$ \\
\bottomrule
\end{tabular}
\end{table*}

\begin{figure*}[t]\centering\pgcfig{figures/exp9_wsweep.pdf}
\caption{Experiment 9: tail constants (left) and the authority margin \(\delta_{\mathrm{cert}}-\delta_{\mathrm{total}}\) of \eqref{eq:control_authority} (right)
versus the cutoff-mode weight \(w\); the dotted line marks \(w^\star\).}\label{fig:exp9}\end{figure*}

\subsection{Experiment 2: Proposition~\ref{prop:twin_crown} (twin-CROWN)}
For \(b_B = {{e2_bB_lo}},\dots,{{e2_bB_hi}}\), with nearest and Adam-learned rounding, the twin-CROWN bound exceeds the
largest Student--Teacher gap over {{e2_nsamp}} samples (uniform samples of \(\mathcal Z_{\mathrm{ODD}}\) plus all box
vertices) by a factor of \({{e2_ratio_lo}}\)--\({{e2_ratio_hi}}\) (Fig.~\ref{fig:exp2}, Table~\ref{tab:exp2_crown}).
The bound is monotone non-increasing in \(b_B\) and halves with every additional bit of \(b_B\). At \(b_B = {{e2_bB}}\), learned rounding
lowers \(\varepsilon_{\mathrm{CROWN}}\) from \({{e2_near}}\) to \({{e2_learn}}\). The port box that indexes the Path-A
table is certified as well.

\begin{table*}[t]
\centering
\caption{Proposition~\ref{prop:twin_crown}: twin-CROWN bound vs.\ the sampled Student--Teacher gaps on $100016$ points of $\mathcal{Z}_{\mathrm{ODD}}$ ($100000$ uniform samples plus all $16$ box vertices).}
\label{tab:exp2_crown}
\footnotesize
\setlength{\tabcolsep}{4pt}
\begin{tabular}{lrrrrrr}
\toprule
$b_B$ & rounding & $\varepsilon_{\mathrm{CROWN}}$ & $E_q$ (explicit) & max sampled gap & mean gap & bound / max gap \\
\midrule
$3$ & nearest & $2.559$ & $2.887$ & $0.9001$ & $0.3945$ & $2.844$ \\
$3$ & learned & $2.326$ & $3.189$ & $0.8927$ & $0.3574$ & $2.606$ \\
$4$ & nearest & $1.329$ & $1.481$ & $0.5366$ & $0.1701$ & $2.476$ \\
$4$ & learned & $1.145$ & $1.603$ & $0.4509$ & $0.1397$ & $2.54$ \\
$5$ & nearest & $0.6217$ & $0.7167$ & $0.2263$ & $0.07976$ & $2.747$ \\
$5$ & learned & $0.5696$ & $0.7491$ & $0.2261$ & $0.0715$ & $2.519$ \\
$6$ & nearest & $0.319$ & $0.3508$ & $0.1449$ & $0.03818$ & $2.201$ \\
$6$ & learned & $0.2868$ & $0.3616$ & $0.1326$ & $0.03538$ & $2.162$ \\
$7$ & nearest & $0.143$ & $0.1801$ & $0.0664$ & $0.01819$ & $2.153$ \\
$7$ & learned & $0.1424$ & $0.1838$ & $0.0586$ & $0.01816$ & $2.43$ \\
$8$ & nearest & $0.07865$ & $0.08686$ & $0.03422$ & $0.00951$ & $2.298$ \\
$8$ & learned & $0.07095$ & $0.09282$ & $0.03621$ & $0.00937$ & $1.96$ \\
$9$ & nearest & $0.03591$ & $0.04118$ & $0.01733$ & $0.004777$ & $2.072$ \\
$9$ & learned & $0.03554$ & $0.04131$ & $0.01776$ & $0.00476$ & $2.001$ \\
$10$ & nearest & $0.01931$ & $0.02345$ & $0.01023$ & $0.002453$ & $1.888$ \\
$10$ & learned & $0.01781$ & $0.02394$ & $0.00809$ & $0.002398$ & $2.202$ \\
$11$ & nearest & $0.009407$ & $0.01112$ & $0.004111$ & $0.001221$ & $2.288$ \\
$11$ & learned & $0.008911$ & $0.01101$ & $0.003928$ & $0.001202$ & $2.268$ \\
$12$ & nearest & $0.004718$ & $0.005478$ & $0.002039$ & $6.126\times10^{-4}$ & $2.314$ \\
$12$ & learned & $0.004455$ & $0.005674$ & $0.001966$ & $6.014\times10^{-4}$ & $2.266$ \\
$13$ & nearest & $0.002385$ & $0.002768$ & $0.001043$ & $3.084\times10^{-4}$ & $2.287$ \\
$13$ & learned & $0.002235$ & $0.002777$ & $0.001134$ & $3.016\times10^{-4}$ & $1.97$ \\
$14$ & nearest & $0.001224$ & $0.001345$ & $5.346\times10^{-4}$ & $1.568\times10^{-4}$ & $2.29$ \\
$14$ & learned & $0.001131$ & $0.001417$ & $5.994\times10^{-4}$ & $1.509\times10^{-4}$ & $1.887$ \\
\bottomrule
\end{tabular}
\end{table*}

\begin{figure}[t]\centering\pgcfig{figures/exp2_crown.pdf}
\caption{Experiment 2: certified \(\varepsilon_{\mathrm{CROWN}}\), explicit \(E_q\) and the largest sampled gap versus \(b_B\).}\label{fig:exp2}\end{figure}

\subsection{Experiment 3: \(\delta_{\mathrm{cert}}\), Lemma~\ref{lem:delta_cert_classH} and Lemma~\ref{lem:delta_cert_closed}}
In class (H), the certified bound of Lemma~\ref{lem:delta_cert_classH} is compared with the brute-force minimum of the
exact bracket of~\eqref{eq:teacher_cbf} over samples of \(\mathcal C_{r,\mathrm{act}}\). The bound is
\({{e3_bound_R2}} \le {{e3_brute_R2}}\) in R2 and \({{e3_bound_R1}} \le {{e3_brute_R1}}\) in R1, so it is sound and nearly
tight. In both regimes \(\delta_{\mathrm{total}} \le \delta_{\mathrm{cert}}\) (Fig.~\ref{fig:exp3}).

On a class-(D) instance, the closed form~\eqref{eq:delta_cert_closed} equals the brute-force infimum of its affine bound
in both branches and lies below the exact-bracket minimum (Table~\ref{tab:exp3_delta_cert}). This is a standalone check
of the lemma, so \(\delta_{\mathrm{total}}\) is reported as N/A there. Using branch 1 unconditionally would overestimate
the margin when \(c < \gamma_0\).

\begin{figure}[t]\centering\pgcfig{figures/exp3_delta_cert.pdf}
\caption{Experiment 3: exact bracket samples versus the certified \(\delta_{\mathrm{cert}}\) and \(\delta_{\mathrm{total}}\) (top);
class-(D) closed form versus the exact bracket (bottom).}\label{fig:exp3}\end{figure}
\begin{table*}[t]
\centering
\caption{Certified margin $\delta_{\mathrm{cert}}$: class-(H) bound of Lemma~\ref{lem:delta_cert_classH} ($c_M=\lambda_{\min}(M,P)-2s_{\mathrm{nl}}$, two-branch infimum) and the class-(D) closed form \eqref{eq:delta_cert_closed} vs.\ the brute-force infimum of the exact bracket over $\mathcal{C}_{r,\mathrm{act}}$. For class (D) the column $c_M$ holds the closed-form slope $c$.}
\label{tab:exp3_delta_cert}
\footnotesize
\setlength{\tabcolsep}{3pt}
\begin{tabular}{lrrrrrrrrr}
\toprule
case & bound & $c_M$ & $\lambda_{\min}(M,P)$ & $s_{\mathrm{nl}}$ & $\gamma_0$ & branch & $\delta_{\mathrm{cert}}$ (bound) & brute-force inf & $\delta_{\mathrm{total}}$ \\
\midrule
R2 & Lemma~\ref{lem:delta_cert_classH} & $0.3725$ & $0.3731$ & $3.005\times10^{-4}$ & $0.35$ & $1$ & $0.01456$ & $0.01468$ & $0.01356$ \\
R1 & Lemma~\ref{lem:delta_cert_classH} & $0.9829$ & $0.9835$ & $3.005\times10^{-4}$ & $0.3$ & $1$ & $0.03161$ & $0.03164$ & $0.002634$ \\
class (D) & \eqref{eq:delta_cert_closed} & $2.59$ & -- & -- & $0.3$ & $1$ & $0.0578$ & $0.06345$ & N/A (standalone) \\
class (D) & \eqref{eq:delta_cert_closed} & $2.59$ & -- & -- & $3$ & $2$ & $0.1032$ & $0.1148$ & N/A (standalone) \\
\bottomrule
\end{tabular}
\end{table*}


\subsection{Experiment 4: Theorem~\ref{thm:main_cbf} (positive case)}
The full method runs with the worst-case adversary on {{n_seeds}} seeds per regime from two kinds of initial
condition: the boundary start \(\tilde h_r(z(0)) = 0\), with the tail at its largest admissible value (H7), and the
interior start \(h_r(z(0)) = {{interior_level}}\,h_{\mathrm{act}}\). Conclusions (a)--(c) hold at every integrator
node for both starts and both regimes (Fig.~\ref{fig:exp4}, Table~\ref{tab:exp4_theorem1}); the largest tail ratio
\(\|x_\perp\|/\epsP\) (norm weighted by \(\Pfe\)) is \({{e4_tailratio_R2}}\) in R2 and \({{e4_tailratio_R1}}\) in R1.
For both starts the minimum of \(h(x(t))\) over the run is attained at the initial instant
(\(t = {{e4_btmin_R2}}\) for the boundary start, \(t = {{e4_itmin_R2}}\) for the interior start): \(\min_t h\) is
\({{e4_bmin_R2}}\) in R2 and \({{e4_bmin_R1}}\) in R1 from the boundary, \({{e4_imin_R2}}\) and \({{e4_imin_R1}}\) from
the interior. The inter-sample dip floor \(-\Delta h_{\mathrm{dip}}\) of conclusion (b) is therefore not exercised
dynamically by these runs. Experiment 4b measures how much of each budget term the same trajectories use.

\begin{figure*}[t]\centering\pgcfig{figures/exp4_theorem1.pdf}
\caption{Experiment 4: raw closed-loop \(h(x(t))\) for the boundary start (top row) and the interior start (middle row),
and the tail ratio \(\|x_\perp\|/\epsP\), with the norm weighted by \(\Pfe\), for the boundary start (bottom row); R2 left, R1 right.}\label{fig:exp4}\end{figure*}
\begin{table*}[t]
\centering
\caption{Theorem~\ref{thm:main_cbf}: measured minima over all seeds and integrator nodes for the boundary start ($\tilde h_r(z(0))=0$) and the interior start ($h_r(z(0))=0.5\,h_{\mathrm{act}}$), worst-case adversary.}
\label{tab:exp4_theorem1}
\footnotesize
\setlength{\tabcolsep}{3pt}
\begin{tabular}{lrrrrrrr}
\toprule
config & start & seeds & $\min_t h(x_t)$ & time of min & $\min \tilde h_r$ & max tail\,/\,$\varepsilon_{\mathrm{basis},P}$ & margin $\delta_{\mathrm{cert}}-\delta_{\mathrm{total}}$ \\
\midrule
R2 & boundary & $16$ & $2.06\times10^{-6}$ & $0$ & $-7.294\times10^{-17}$ & $1$ & $9.994\times10^{-4}$ \\
R2 & interior & $16$ & $0.01034$ & $0$ & $6.709\times10^{-4}$ & $0.0583$ & $9.994\times10^{-4}$ \\
R1 & boundary & $16$ & $3.008\times10^{-6}$ & $0$ & $-7.774\times10^{-17}$ & $1$ & $0.02898$ \\
R1 & interior & $16$ & $0.0128$ & $0$ & $0.01274$ & $0.04749$ & $0.02898$ \\
\bottomrule
\end{tabular}
\end{table*}


\subsection{Experiment 4b: budget utilisation along the certified trajectories}
Experiment 4b re-runs the closed loops of Experiment 4 and evaluates, at each of the {{e4b_nodes_R2}} (R2) and
{{e4b_nodes_R1}} (R1) integrator nodes over both starts and all seeds, every term of~\eqref{eq:ode_diff_ineq} next to
the bound the proof uses for it (Table~\ref{tab:exp4b_budget}). Each fraction below is the realised term over its
bound: the maximum over all nodes and, in parentheses, the median over the {{e4b_act_R2}} (R2) and {{e4b_act_R1}} (R1)
nodes whose sample \(z_k\) lies in \(\mathcal C_{r,\mathrm{act}}\), where the budget enters the proof. A fraction
{{e4b_intfrac_R2}} (R2) and {{e4b_intfrac_R1}} (R1) of the nodes belong to periods with \(z_k \in \mathcal B_{\mathrm{int}}\);
there the proof uses only the Lipschitz branch, and its inequality holds at every such node.
\begin{itemize}
\item Adversary: \({{e4b_adv_max_R2}}\) (\({{e4b_adv_med_R2}}\)) in R2 and \({{e4b_adv_max_R1}}\) (\({{e4b_adv_med_R1}}\))
in R1. The worst-case \(w\) has norm \(\varepsilon_{\mathrm{adv}}\) at every node, so the gap to one is the supremum
of \(\|d_r^\top\nabla\tilde h_r\|_2\) over \(\Crz\) in \(\overline d_h\).
\item Quantization and lag: \({{e4b_quant_max_R2}}\) (\({{e4b_quant_med_R2}}\)) and \({{e4b_quant_max_R1}}\)
(\({{e4b_quant_med_R1}}\)). The mismatch \(e_u\) reaches at most \({{e4b_eu_max_R2}}\) (R2) and \({{e4b_eu_max_R1}}\) (R1)
of its phase bound, so \(\varepsilon_{\mathrm{CROWN}}\), \(\varepsilon_A\) and \(\Delta u_{\mathrm{lag}}\) are the conservative constants of this term.
\item Tail coupling: \({{e4b_tail_max_R2}}\) (\({{e4b_tail_med_R2}}\)) and \({{e4b_tail_max_R1}}\) (\({{e4b_tail_med_R1}}\)).
Of the factors of \(L_h\Lperp\epsP\), the two least used are the Lipschitz constant \(\Lperp\) (median
\(\|\Delta_\perp\|_2/(\Lperp\|x_\perp\|) = {{e4b_lip_med_R2}}\) in R2 and \({{e4b_lip_med_R1}}\) in R1) and the
alignment of \(\nabla\tilde h_r\) with \(\Delta_\perp\) (median cosine \({{e4b_cos_med_R2}}\) and \({{e4b_cos_med_R1}}\)).
\item Hold shift: \({{e4b_hold_max_R2}}\) (\({{e4b_hold_med_R2}}\)) and \({{e4b_hold_max_R1}}\) (\({{e4b_hold_med_R1}}\)), the
smallest of the four fractions in both regimes. Two constants of \(L_{\mathrm{hold}}\) are conservative here:
\(\|z - z_k\|_2\) reaches at most \({{e4b_mz_max_R2}}\) (R2) and \({{e4b_mz_max_R1}}\) (R1) of \(M_z s\), and the hold shift
at most \({{e4b_lipH_max_R2}}\) (R2) and \({{e4b_lipH_max_R1}}\) (R1) of
\(L_{\mathrm{hold}}\|z - z_k\|_2/M_z\), its Lipschitz bound in \(z\).
\end{itemize}
The realised rate \(\dot{\tilde h}_r + \gamma_0\tilde h_r\) minus the right-hand side of~\eqref{eq:ode_diff_ineq} is
positive at every active-shell node, so the soundness check passes. Its minimum is \({{e4b_minslack_R2}}\,\delta_{\mathrm{total}}\)
in R2 and \({{e4b_minslack_R1}}\,\delta_{\mathrm{total}}\) in R1, attained at a sampling instant, where the hold term
vanishes; the Teacher margin at \(z_k\) (the left-hand side of~\eqref{eq:teacher_cbf} minus \(\delta_{\mathrm{total}}\))
supplies a fraction \({{e4b_teachshare_R2}}\) (R2) and \({{e4b_teachshare_R1}}\) (R1) of it.

\begin{table*}[t]
\centering
\caption{Experiment 4b: realised fraction of each budget term of \eqref{eq:ode_diff_ineq} along the certified trajectories of Experiment 4 (both starts, all seeds, every integrator node, worst-case adversary); row (v): minimum and median of the realised rate minus the right-hand side of \eqref{eq:ode_diff_ineq}, in units of $\delta_{\mathrm{total}}$. Columns $\mathcal{C}_{r,\mathrm{act}}$: nodes of periods whose sample $z_k$ lies in the active shell, where the budget enters the proof.}
\label{tab:exp4b_budget}
\footnotesize
\setlength{\tabcolsep}{4pt}
\begin{tabular}{lrrrrr}
\toprule
config & term & max, all nodes & median, all nodes & max, $\mathcal{C}_{r,\mathrm{act}}$ & median, $\mathcal{C}_{r,\mathrm{act}}$ \\
\midrule
R2 & (i) tail coupling & $0.3161$ & $2.063\times10^{-4}$ & $0.3161$ & $0.003444$ \\
R2 & (ii) adversary & $0.8772$ & $0.01896$ & $0.8772$ & $0.1454$ \\
R2 & (iii) quantization/lag & $0.3729$ & $0.001841$ & $0.3729$ & $0.01712$ \\
R2 & (iv) hold shift & $0.07439$ & $5.455\times10^{-6}$ & $0.07439$ & $0.001438$ \\
R2 & (iii) phase $[0,\tau_A)$ & $0.2057$ & $7.39\times10^{-4}$ & $0.2057$ & $0.008388$ \\
R2 & (iii) phase $[\tau_A,\tau_B)$ & $0.3729$ & $0.003766$ & $0.3729$ & $0.0348$ \\
R2 & (iii) phase $[\tau_B,T_s)$ & $0.2551$ & $0.002002$ & $0.2551$ & $0.01809$ \\
R2 & (v) slack$\,/\,\delta_{\mathrm{total}}$: min, median & N/A & N/A & $0.6597$ & $2.807$ \\
R1 & (i) tail coupling & $0.2606$ & $4.582\times10^{-4}$ & $0.2606$ & $0.01455$ \\
R1 & (ii) adversary & $0.9761$ & $0.04194$ & $0.9761$ & $0.2814$ \\
R1 & (iii) quantization/lag & $0.3672$ & $0.005189$ & $0.3672$ & $0.03416$ \\
R1 & (iv) hold shift & $0.09728$ & $8.842\times10^{-4}$ & $0.09728$ & $0.0226$ \\
R1 & (iii) phase $[0,\tau_A)$ & $0.2559$ & $0.002249$ & $0.2559$ & $0.02988$ \\
R1 & (iii) phase $[\tau_A,\tau_B)$ & $0.3672$ & $0.009477$ & $0.3672$ & $0.05098$ \\
R1 & (iii) phase $[\tau_B,T_s)$ & $0.2856$ & $0.005166$ & $0.2856$ & $0.02683$ \\
R1 & (v) slack$\,/\,\delta_{\mathrm{total}}$: min, median & N/A & N/A & $11.5$ & $17.56$ \\
\bottomrule
\end{tabular}
\end{table*}


\subsection{Experiment 5: Proposition~\ref{prop:deadzone_stability} and ablation (d)}
In R1, (H-dz) holds and the trajectories enter \(\mathcal B_{\mathrm{ss}}\). With the mid-tread deadzone, Path A is
silent late in the run (\(u_A \equiv 0\)) for every \(b_A \in \{{{e5_bAs}}\}\). Ablation (d) replaces it with a
mid-rise quantiser and gives a persistent actuator limit cycle at every \(b_A\), including the deployed
\(b_A = {{e5_bA_dep}}\) (Figs.~\ref{fig:exp5} and~\ref{fig:exp5z}, Table~\ref{tab:exp5_deadzone}). The number of sign switches per unit time
falls from \({{e5_sw_lo}}\) at \(b_A = {{e5_bA_lo}}\) to \({{e5_sw_hi}}\) at \(b_A = {{e5_bA_hi}}\), and the final lag error \(\|u_A - \hat{\mathcal T}\|_2\) of the mid-rise quantiser at
\(b_A = {{e5_bA_hi}}\) is \({{e5_lag_dep}}\). With \(b_A = 1\), \(r_{\mathrm{ss}} = {{e5_rss1}} > r_{\min} = {{e5_rmin1}}\), so (H-dz) fails and
Proposition~\ref{prop:deadzone_stability} is refused.

R1 is open-loop stable, so the chattering appears in the actuator signal, not as a larger state amplitude.

\begin{figure*}[t]\centering\pgcfig{figures/exp5_deadzone_sweep.pdf}
\caption{Experiment 5 / ablation (d): late sign switches of the mid-rise Path-A command (left) and the late-window state
size (right) versus \(b_A\).}\label{fig:exp5}\end{figure*}
\begin{figure*}[t]\centering\pgcfig{figures/exp5_deadzone_z.pdf}
\caption{Experiment 5: \(\|z(t_k)\|_P\) for the deadzone (left) and mid-rise (right) quantiser over the \(b_A\) sweep.}\label{fig:exp5z}\end{figure*}
\begin{table*}[t]
\centering
\caption{Proposition~\ref{prop:deadzone_stability} and ablation (d) on R1 (disturbance $w\equiv0$): mid-tread LUT with deadzone vs.\ mid-rise LUT (no zero level) at the deployed $b_A=9$ and over the sweep $b_A\in\{3,5,7,9\}$; the last row is $b_A=1$, which violates (H-dz). Final lag: $\min$ over seeds of $\|u_A(t_k)-\hat{\mathcal{T}}(z(t_k))\|_2$ at the last sample. For every swept $b_A$, $u_A\equiv0$ in the late window with the deadzone and $u_A\neq0$ at every late sample without it (asserted).}
\label{tab:exp5_deadzone}
\footnotesize
\setlength{\tabcolsep}{3pt}
\begin{tabular}{lrrrrrrrrr}
\toprule
LUT & $b_A$ & $\varepsilon_A$ & $r_{\mathrm{ss}}$ & $r_{\min}$ & (H-dz) & late $\max\|z\|_P$ & final lag $\|u_A-\hat{\mathcal{T}}\|_2$ & switches/time & $\min h$ \\
\midrule
deadzone & $3$ & $0.02022$ & $0.105$ & $0.156$ & \checkmark & $2.035\times10^{-4}$ & $1.966\times10^{-6}$ & $0$ & $3.939\times10^{-6}$ \\
deadzone & $5$ & $0.005395$ & $0.03145$ & $0.156$ & \checkmark & $2.011\times10^{-4}$ & $1.914\times10^{-6}$ & $0$ & $3.939\times10^{-6}$ \\
deadzone & $7$ & $0.001689$ & $0.01306$ & $0.156$ & \checkmark & $1.974\times10^{-4}$ & $1.854\times10^{-6}$ & $0$ & $3.939\times10^{-6}$ \\
deadzone & $9$ & $7.624\times10^{-4}$ & $0.01162$ & $0.156$ & \checkmark & $1.934\times10^{-4}$ & $1.833\times10^{-6}$ & $0$ & $3.939\times10^{-6}$ \\
mid-rise & $3$ & $0.02022$ & $0.105$ & $0.156$ & \checkmark & $1.257\times10^{-4}$ & $0.0197$ & $953.1$ & $3.939\times10^{-6}$ \\
mid-rise & $5$ & $0.005395$ & $0.03145$ & $0.156$ & \checkmark & $8.268\times10^{-5}$ & $0.004851$ & $859.9$ & $3.939\times10^{-6}$ \\
mid-rise & $7$ & $0.001689$ & $0.01306$ & $0.156$ & \checkmark & $1.559\times10^{-4}$ & $0.001286$ & $469.4$ & $3.939\times10^{-6}$ \\
mid-rise & $9$ & $7.624\times10^{-4}$ & $0.01162$ & $0.156$ & \checkmark & $1.845\times10^{-4}$ & $3.507\times10^{-4}$ & $121.6$ & $3.939\times10^{-6}$ \\
deadzone & $1$ & $0.07951$ & $0.3992$ & $0.156$ & $\times$ & $2.067\times10^{-4}$ & $2.005\times10^{-6}$ & $0$ & $3.939\times10^{-6}$ \\
\bottomrule
\end{tabular}
\end{table*}


\subsection{Experiment 6: ablations (a), (b), (c), (e)}
Each ablation removes one component at the deployed \(w^\star\). In every case the full certificate refuses and the
full-order plant leaves \(\mathcal C\) (Fig.~\ref{fig:exp6}, Table~\ref{tab:exp6_ablations}). For (a) and (b), a naive
certificate with the corresponding term dropped would accept.
\begin{itemize}
\item (a) \(K_m = {{abl_K_m}}\) without tail bookkeeping: {{e6_a_exits}} of {{n_seeds}} seeds exit.
\item (b) \(b_B = {{abl_b_B}}\) without the CROWN margin: {{e6_b_exits}} of {{n_seeds}} seeds exit.
\item (c) Single-rate hold (Table~\ref{tab:exp6_multirate}). At the deployed \(T_s\), \(\delta_{\mathrm{total}}\) rises from
\({{e6_c_multi}}\) to \({{e6_c_single}} > \delta_{\mathrm{cert}}\), so (H6) fails. At \(T_s = {{e6_TsV}}\), single rate gives
{{e6_c_exits}} exits and multi-rate gives none.
\item (e) \(\epsilon_{\mathrm{adv}} = {{abl_eps_adv}}\): {{e6_e_exits}} of {{n_seeds}} seeds exit. At the largest certifiable
\(\epsilon_{\mathrm{adv}}^\star = {{e6_eps_star}}\) there are no exits.
\end{itemize}

\begin{figure*}[t]\centering\pgcfig{figures/exp6_ablations.pdf}
\caption{Experiment 6: raw \(h(x(t))\) under each ablation; violations appear below \(h=0\).}\label{fig:exp6}\end{figure*}
\begin{table*}[t]
\centering
\caption{Ablations at the deployed $w^\star$: every removed component yields a certificate refusal and a full-order safety violation under the worst-case adversary. Row (e): the largest certifiable adversary level is $\epsilon_{\rm adv}^\star=0.004223$.}
\label{tab:exp6_ablations}
\footnotesize
\setlength{\tabcolsep}{3pt}
\begin{tabular}{lrrrrrrrr}
\toprule
ablation & refused & failing & $\delta_{\mathrm{total}}$ & $\delta_{\mathrm{cert}}$ & naive $\delta_{\rm total}$ & naive $\delta_{\rm cert}$ & exits (of $16$) & $\min_t h$ \\
\midrule
(a) $K_m=1$, no tail bookkeeping & \checkmark & (H3), (H4), (H6), (H7) & $7.593$ & $-60.63$ & $0.007956$ & $0.0335$ & $4$ & $-0.02545$ \\
(b) $b_B=2$, no CROWN margin & \checkmark & (H6) & $2.149$ & $0.01456$ & $0.01135$ & $0.01456$ & $9$ & $-0.04155$ \\
(c) single rate $\tau_A=\tau_B=T_s$ & \checkmark & (H6) & $0.01702$ & $0.01456$ & -- & -- & $10$ & $-0.1591$ \\
(e) $\epsilon_{\rm adv}=1$ & \checkmark & (H6) & $0.4617$ & $0.01409$ & -- & -- & $16$ & $-0.02605$ \\
\bottomrule
\end{tabular}
\end{table*}

\begin{table*}[t]
\centering
\caption{Ablation (c): exponential weights $w_1,w_2$, $\Delta h_{\rm dip}$ and $\delta_{\rm total}$ for multi-rate vs.\ single-rate hold.}
\label{tab:exp6_multirate}
\footnotesize
\setlength{\tabcolsep}{3pt}
\begin{tabular}{lrrrrrrrr}
\toprule
schedule & $\tau_A$ & $\tau_B$ & $w_1$ & $w_2$ & $\Delta h_{\mathrm{dip}}$ & $\delta_{\mathrm{total}}$: lag term & $\delta_{\mathrm{total}}$ & (H6) \\
\midrule
multi rate & $1.5\times10^{-5}$ & $3\times10^{-5}$ & $0.05$ & $0.05$ & $2.06\times10^{-6}$ & $1.859\times10^{-4}$ & $0.01356$ & \checkmark \\
single rate & $3\times10^{-4}$ & $3\times10^{-4}$ & $1$ & $0$ & $3.097\times10^{-6}$ & $0.003718$ & $0.01702$ & $\times$ \\
single rate, $T_s=0.2$ & $0.2$ & $0.2$ & $1$ & $0$ & $1.377$ & $2.479$ & $6.941$ & $\times$ \\
multi rate, $T_s=0.2$ & $0.01$ & $0.02$ & $0.04835$ & $0.04852$ & $0.9057$ & $0.1199$ & $4.582$ & $\times$ \\
\bottomrule
\end{tabular}
\end{table*}


\subsection{Experiment 7: Corollary~\ref{cor:stochastic_w2} and ablation (f)}
In R1 the paths start on the level \({{e7_z0}}\) of \(\tilde{\mathcal C}_r\) (\(h(x_0) = {{e7_h0}}\)), with spatially smooth
noise. The intensity is chosen by a grid search over \(\sigma \in \{{{e7_grid}}\}\) (Table~\ref{tab:exp7_sigma_grid}).
\(\sigma^\star = {{e7_sigma}}\) is the largest grid value for which the Corollary accepts at \(T = {{e7_T}}\) and
\(\max_t|h(x_t) - h(x_0)| \ge {{e7_ffrac}}\,h(x_0)\) holds on at least half of the paths; the Corollary refuses from
\(\sigma = {{e7_first_refused}}\) on (failing condition: {{e7_refused_cond}}). On the screening paths at \(\sigma^\star\) the criterion is driven mostly by the closed-loop drift (median
deviation \({{e7_meddev}}\,h(x_0)\)). The noise itself is not negligible: the median deviation from the noise-free
trajectory is \({{e7_noisedev}}\,h(x_0)\) (Table~\ref{tab:exp7_sigma_grid}).

Out of {{e7_paths}} Monte Carlo paths, {{e7_exits}} exit. The exact Clopper--Pearson \({{e7_cp_pct}}\%\) upper limit is
\({{e7_cp_hi}} \le \alpha = {{e7_alpha}}\) (Fig.~\ref{fig:exp7}, Table~\ref{tab:exp7_stochastic}). The bound is very
conservative: \(T_{\mathrm{TTL}} = {{e7_TTL}}\) exceeds the horizon by a factor of \({{e7_TTL_ratio}}\), and no exit is
observed.

Ablation (f) uses spatially white small-scale noise. The full Corollary refuses it: the tail envelope
\(\Delta_\sigma = {{e7f_Dsig}}\) exceeds \(\sqrt{E_{\mathrm{safe}}} = {{e7f_sqrtE}}\). A cruder bound that ignores the tail
diffusion (\(\bar\sigma_\perp := 0\)) accepts it, yet {{e7f_exits}} of {{e7_paths}} paths exit (Clopper--Pearson lower
limit \({{e7f_cp_lo}} > \alpha\)). The tail-diffusion term is therefore load-bearing. The constants are split into
Tables~\ref{tab:exp7_corollary_det} and~\ref{tab:exp7_corollary_stoch}.

The Corollary is not exercised on R2. There the pointwise penalty \(\gamma_0h_{\mathrm{act}} = {{e7_R2_gamma0_h_act}}\)
already exceeds the margin \(\delta_{\mathrm{cert}} - \delta_{\mathrm{total}} = {{e7_R2_margin}}\), so its authority condition
fails even at \(\sigma = 0\) (asserted; Table~\ref{tab:exp1_constants_b}). All stochastic runs use \(r_{W_{\scriptstyle 2}} = {{r_W2}}\),
i.e.\ no distributional drift of the adversary.

\begin{figure*}[t]\centering\pgcfig{figures/exp7_stochastic.pdf}
\caption{Experiment 7: raw stochastic \(h(x_t)\) paths for profile noise at \(\sigma^\star\) with the noise-free run (left)
and for ablation (f) (right).}\label{fig:exp7}\end{figure*}
\begin{table*}[t]
\centering
\caption{Grid search over the profile-noise intensity $\sigma$ (R1, $T=2$, $\alpha=0.05$). Fluct.\ share: fraction of 100 screening paths with $\max_t|h(x_t)-h(x_0)|\ge 0.1\,h(x_0)$; noise dev.: median of $\max_t|h(x_t)-h_{\rm det}(t)|/h(x_0)$ (deviation from the noise-free run). $\sigma^\star=0.007$.}
\label{tab:exp7_sigma_grid}
\footnotesize
\setlength{\tabcolsep}{3pt}
\begin{tabular}{lrrrrrrrr}
\toprule
$\sigma$ & Corollary & failing condition & $\delta_{\rm total,\sigma}+\delta_{\rm stoch}+\delta_{\rm drift}$ & $\delta_{\mathrm{cert},\sigma}$ & $T_{\mathrm{TTL}}$ & fluct. share & total dev. & noise dev. \\
\midrule
$0.002$ & \checkmark & none & $0.0161$ & $0.03161$ & $1015$ & $1$ & $8.399$ & $0.04911$ \\
$0.003$ & \checkmark & none & $0.0175$ & $0.03161$ & $1248$ & $1$ & $8.401$ & $0.07428$ \\
$0.005$ & \checkmark & none & $0.02084$ & $0.03161$ & $1667$ & $1$ & $8.405$ & $0.1269$ \\
$0.007$ & \checkmark & none & $0.02496$ & $0.03161$ & $2022$ & $1$ & $8.41$ & $0.174$ \\
$0.01$ & $\times$ & margin & $0.03276$ & $0.03161$ & $2440$ & -- & -- & -- \\
$0.014$ & $\times$ & margin & $0.04671$ & $0.03161$ & $2793$ & -- & -- & -- \\
$0.02$ & $\times$ & margin & $0.07733$ & $0.03161$ & $2933$ & -- & -- & -- \\
\bottomrule
\end{tabular}
\end{table*}

\begin{table*}[t]
\centering
\caption{Corollary~\ref{cor:stochastic_w2} (R1, $T=2$, $\alpha=0.05$, $r_{W_{\scriptstyle 2}}=0$, worst-case adversary, $z(0)$ on the level $0.9$ of $\tilde{\mathcal{C}}_r$): Monte Carlo exit frequency vs.\ the bound $\alpha$ with the exact Clopper--Pearson interval; crude: bound with $\sigma_\perp:=0$. On R2 the Corollary refuses already at $\sigma=0$: $\gamma_0h_{\mathrm{act}}=0.007237$ exceeds the margin $\delta_{\mathrm{cert}}-\delta_{\mathrm{total}}=9.994\times10^{-4}$.}
\label{tab:exp7_stochastic}
\footnotesize
\setlength{\tabcolsep}{3pt}
\begin{tabular}{lrrrrrrrr}
\toprule
noise & $\sigma$ & $h(x_0)$ & Corollary & crude & $T_{\mathrm{TTL}}$ & $T_{\mathrm{TTL}}/T$ & exits & CP $95\%$ interval \\
\midrule
profile, $\sigma^\star$ & $0.007$ & $0.005051$ & \checkmark & \checkmark & $2022$ & $1011$ & $0/2000$ & $[0,\,0.001843]$ \\
tail-white (f) & $0.02$ & $0.005051$ & $\times$ & \checkmark & $0$ & $0$ & $2000/2000$ & $[0.9982,\,1]$ \\
\bottomrule
\end{tabular}
\end{table*}

\begin{table*}[t]
\centering
\caption{Corollary~\ref{cor:stochastic_w2}, diffusion-enlarged deterministic constants (full vs.\ crude bound); the tail envelope must stay below $\sqrt{E_{\mathrm{safe}}}=0.2236$.}
\label{tab:exp7_corollary_det}
\footnotesize
\setlength{\tabcolsep}{4pt}
\begin{tabular}{lrrrrrr}
\toprule
bound & $\sigma_r$ & $\sigma_\perp^2$ & $\Delta_\sigma$ & $\varepsilon_{\mathrm{basis},\sigma}$ & $\delta_{\mathrm{total},\sigma}$ & $T_\perp$ \\
\midrule
profile, $\sigma^\star$, full & $0.007$ & $4.917\times10^{-10}$ & $9.703\times10^{-5}$ & $0.00739$ & $0.002634$ & $\infty$ \\
profile, $\sigma^\star$, crude $\sigma_\perp:=0$ & $0.007$ & $0$ & $0$ & $0.007293$ & $0.002634$ & $\infty$ \\
tail-white (f), full & $8.742\times10^{-18}$ & $0.0632$ & $1.1$ & $1.107$ & $0.005408$ & $0$ \\
tail-white (f), crude $\sigma_\perp:=0$ & $8.742\times10^{-18}$ & $0$ & $0$ & $0.007293$ & $0.002634$ & $\infty$ \\
\bottomrule
\end{tabular}
\end{table*}

\begin{table*}[t]
\centering
\caption{Corollary~\ref{cor:stochastic_w2}, stochastic constants of \eqref{eq:delta_stoch} and the Kushner bound (full vs.\ crude bound).}
\label{tab:exp7_corollary_stoch}
\footnotesize
\setlength{\tabcolsep}{4pt}
\begin{tabular}{lrrrrrr}
\toprule
bound & $r_{\mathrm{BM}}$ & $\delta_{\mathrm{It\hat o}}$ & $\sigma_h$ & $\Delta\delta_{\mathrm{pw}}$ & $\delta_{\mathrm{stoch}}$ & $c_{\mathrm{int}}$ \\
\midrule
profile, $\sigma^\star$, full & $0.003237$ & $4.9\times10^{-5}$ & $0.00313$ & $0.008873$ & $0.02074$ & $5.146\times10^{-6}$ \\
profile, $\sigma^\star$, crude $\sigma_\perp:=0$ & $0.003237$ & $4.9\times10^{-5}$ & $0.00313$ & $0.008873$ & $0.02074$ & $5.146\times10^{-6}$ \\
tail-white (f), full & $4.043\times10^{-18}$ & $8.786\times10^{-35}$ & $3.909\times10^{-18}$ & $0.001904$ & $-0.1744$ & $-1.837\times10^{-7}$ \\
tail-white (f), crude $\sigma_\perp:=0$ & $4.043\times10^{-18}$ & $8.786\times10^{-35}$ & $3.909\times10^{-18}$ & $0.001896$ & $0.009573$ & $4.578\times10^{-6}$ \\
\bottomrule
\end{tabular}
\end{table*}


\subsection{Experiment 8: numerical method}
At the convergence-test step, explicit Euler is unstable (amplification \({{e8_amp}}\)). The stiffness is moderate
(spectral radius \({{e8_rad}}\), stiffness ratio \({{e8_stiff}}\)). Step halving gives error ratios of
\({{e8_r_lo}}\)--\({{e8_r_hi}}\) for ETD2RK, i.e.\ order 2 (Fig.~\ref{fig:exp8}, Table~\ref{tab:exp8_convergence}).
The simulator matches the closed-form augmented exponential of a latched linear system to \({{e8_latch}}\), whereas
rounding the latch instants to a step grid gives \({{e8_grid}}\). The sparse modal propagation equals a dense nodal
ETD2RK to \({{e8_modal}}\), and two runs give identical SHA-256 hashes (Table~\ref{tab:exp8_numerics}).

\begin{figure}[t]\centering\pgcfig{figures/exp8_convergence.pdf}
\caption{Experiment 8: step-halving convergence of ETD2RK and ETD1.}\label{fig:exp8}\end{figure}
\begin{table*}[t]
\centering
\caption{Step-halving convergence of the exponential integrator (open-loop nonlinear lattice).}
\label{tab:exp8_convergence}
\footnotesize
\setlength{\tabcolsep}{4pt}
\begin{tabular}{lrrrr}
\toprule
substeps per period & step & rel. error ETD2RK & ratio & rel. error ETD1 \\
\midrule
$1$ & $0.05$ & $1.238\times10^{-9}$ & -- & $4.081\times10^{-8}$ \\
$2$ & $0.025$ & $3.095\times10^{-10}$ & $4.002$ & $2.048\times10^{-8}$ \\
$4$ & $0.0125$ & $7.731\times10^{-11}$ & $4.003$ & $1.026\times10^{-8}$ \\
$8$ & $0.00625$ & $1.927\times10^{-11}$ & $4.011$ & $5.135\times10^{-9}$ \\
$16$ & $0.003125$ & $4.765\times10^{-12}$ & $4.045$ & $2.569\times10^{-9}$ \\
\bottomrule
\end{tabular}
\end{table*}

\begin{table}[t]
\centering
\caption{Numerical-integrity checks.}
\label{tab:exp8_numerics}
\footnotesize
\setlength{\tabcolsep}{4pt}
\begin{tabular}{lr}
\toprule
check & value \\
\midrule
spectral radius $\max|\lambda|$ & $72$ \\
explicit Euler amplification at conv. step & $3.73$ \\
stiffness ratio & $43.83$ \\
latch: rel. error exact & $2.808\times10^{-14}$ \\
latch: rel. error grid-rounded & $0.01675$ \\
modal vs nodal rel. error & $1.928\times10^{-14}$ \\
ETD2RK observed orders & $2.001,\ 2.001,\ 2.004,\ 2.016$ \\
determinism SHA-256 (prefix) & \texttt{2bde49483745ec98} \\
\bottomrule
\end{tabular}
\end{table}

\begin{table}[t]
\centering
\caption{Measured wall-clock and peak resident memory per script (2 cores); experiment assertions passed: $248$ of $248$; pytest: $65$ passed.}
\label{tab:runtimes}
\footnotesize
\begin{tabular}{lrr}
\toprule
script & wall [s] & peak RSS [MB] \\
\midrule
exp0 openloop & $18.4$ & $214$ \\
exp1 constants & $4.8$ & $112$ \\
exp2 crown & $6.9$ & $150$ \\
exp3 delta cert & $5.9$ & $206$ \\
exp4 theorem1 & $52.6$ & $343$ \\
exp4b budget & $203.6$ & $1010$ \\
exp5 deadzone & $54.4$ & $210$ \\
exp6 ablations & $65.2$ & $299$ \\
exp7 stochastic & $166.4$ & $1040$ \\
exp8 numerics & $9.2$ & $220$ \\
exp9 wsweep & $31.4$ & $194$ \\
pytest & $23.7$ & $1193$ \\
\midrule
total & $642.9$ & $1193$ \\
\bottomrule
\end{tabular}
\end{table}

